\documentclass[11pt]{article}
\usepackage[utf8]{inputenc}
\usepackage[margin=1in]{geometry}
\usepackage{amsmath}
\usepackage{graphicx}
\usepackage{booktabs}
\usepackage{hyperref}
\usepackage[round]{natbib}

\title{Chronic Chemogenetic Activation of Negative Memory Engrams in the Ventral Hippocampus Produces Lasting Anxiety, Cognitive Deficits, and Neuroinflammatory Changes in Mice}
\author{Author A,$^{1}$ Author B,$^{1}$ Author C,$^{2}$ Author D$^{1,2,*}$\\
\small $^{1}$Department of Neuroscience, University of Neuroscience\\
\small $^{2}$Center for Memory and Brain, University of Neuroscience\\
\small $^{*}$Corresponding author: authord@university.edu}
\date{}

\begin{document}
\maketitle

\begin{abstract}
Chronic psychological stress is a major risk factor for neuropsychiatric disorders including major depressive disorder, generalized anxiety disorder, and post-traumatic stress disorder. While external stressors have been extensively studied, whether the repeated internal reactivation of a negative memory alone can recapitulate chronic-stress-like pathology remains unknown. Here we used TRAP2 transgenic mice to permanently label neurons active during contextual fear conditioning with the excitatory DREADD hM3Dq in the ventral hippocampus (vHPC), enabling chronic chemogenetic reactivation of a negative memory engram over three months. We first show via RNA-sequencing and gene set enrichment analysis that negative vHPC engram cells are enriched for pro-inflammatory gene expression (including \textit{Spp1}, \textit{Ttr}, and \textit{C1qb}), while positive engram cells are enriched for anti-inflammatory signatures. Following chronic reactivation, both young (6-month) and aged (14-month) mice exhibited heightened anxiety-like behavior in the open field test and zero maze, as well as fear generalization to novel contexts. Notably, spatial working memory impairment in the Y-maze was observed only in 14-month-old mice, and fear extinction deficits were specific to 6-month-old mice, revealing age-dependent vulnerability profiles. Immunohistochemical analysis demonstrated significant alterations in Iba1$^+$ microglial and GFAP$^+$ astrocytic number and morphology in hippocampal subregions, consistent with neuroinflammation, without detectable neuronal loss. These findings demonstrate that chronic reactivation of a single negative memory engram is sufficient to produce a multifaceted behavioral and cellular pathology resembling chronic stress, with distinct age-dependent manifestations.
\end{abstract}

\section{Introduction}

Chronic stress constitutes one of the most pervasive risk factors for the development of neuropsychiatric disorders, including major depressive disorder (MDD), generalized anxiety disorder (GAD), and post-traumatic stress disorder (PTSD) \citep{kendler2001causal, mcewen2007physiology}. Decades of research have established that prolonged exposure to external stressors produces lasting changes in brain circuitry, neuroendocrine function, and immune signaling. However, a fundamental question remains: can the repeated internal reactivation of a negative memory, in the absence of any external stressor, recapitulate the behavioral and cellular consequences of chronic stress?

The discovery of memory engrams---ensembles of neurons whose activity patterns encode specific memories---has transformed our understanding of how experiences are stored and recalled in the brain \citep{josselyn2015memory}. Seminal work demonstrated that artificial reactivation of engram cells in the hippocampus can elicit behavioral responses associated with the original memory \citep{ramirez2013creating}, and that the emotional valence of engram-associated memories can be bidirectionally modulated \citep{redondo2014bidirectional}. Furthermore, studies have shown that hippocampal memory traces are differentially modulated by experience, time, and neurogenesis \citep{denny2014hippocampal}. Yet while acute reactivation of negative engrams has been shown to modulate social avoidance and fear responses, no study has examined the consequences of sustained, chronic stimulation of a negative memory engram over extended periods.

The ventral hippocampus (vHPC) is of particular interest in this context, as it serves as a critical node integrating memory, emotion, and stress responses \citep{fanselow2000contextual}. The vHPC sends dense projections to the amygdala, prefrontal cortex, and hypothalamus, positioning it to influence both behavioral and neuroendocrine stress responses. Importantly, positive and negative memory engrams in the vHPC have been shown to be molecularly distinct, yet the specific inflammatory gene signatures associated with these engrams have not been characterized comprehensively.

An additional gap in the literature concerns the role of aging. The aged brain responds differently to stress, with altered microglial reactivity, reduced neurogenesis, and changes in synaptic plasticity \citep{roth2009lasting}. Whether chronic reactivation of a negative engram produces different effects in young versus aged animals has not been tested. Understanding such age-dependent vulnerability could have important implications for age-related neuropsychiatric conditions.

Glial cells, particularly microglia and astrocytes, are increasingly recognized as key mediators of stress-related brain pathology \citep{nimmerjahn2005resting, wohleb2016neuron}. Microglia, the brain's resident immune cells, undergo morphological and functional changes in response to chronic stress, shifting from a ramified surveillance state to a more activated, inflammatory phenotype \citep{york2021microglial}. Astrocytes similarly respond to chronic stress with changes in morphology and GFAP expression. However, the glial consequences of chronic memory reactivation---as opposed to chronic external stress---remain entirely unexplored.

In this study, we employed TRAP2 (Targeted Recombination in Active Populations 2) transgenic mice \citep{allen2017trap2} combined with Cre-dependent expression of the excitatory DREADD hM3Dq \citep{armbruster2007evolving} to permanently tag and chronically reactivate neurons encoding a negative memory in the vHPC. Using RNA-sequencing with gene set enrichment analysis, a comprehensive behavioral battery, and detailed immunohistochemical analysis of glial morphology with 3DMorph-based quantification \citep{york2018microglia}, we tested whether chronic negative-engram reactivation produces lasting behavioral and cellular abnormalities, and whether these effects differ between young and aged animals. Our results reveal that chronic reactivation of a single negative memory engram is sufficient to induce an anxiety-like phenotype, impair cognitive function in an age-dependent manner, and drive neuroinflammatory glial changes---all without detectable neuronal loss.

\section{Methods}

\subsection{Animals}

TRAP2 transgenic mice on a C57BL/6J background were used for all behavioral and immunohistochemical experiments. A separate cohort of wild-type C57BL/6J mice was used for RNA-sequencing experiments. All mice were maintained on a 12-hour light/dark cycle with ad libitum access to food and water. All procedures were approved by the Institutional Animal Care and Use Committee.

\subsection{Viral Injections}

Mice received bilateral stereotaxic injections of AAV9-hSyn-DIO-hM3Dq-mCherry (experimental group) or AAV9-hSyn-DIO-mCherry (control group) targeting the ventral hippocampus. Injection coordinates were determined relative to bregma using a standard mouse brain atlas. A volume of 0.5~$\mu$L was injected per hemisphere at a rate of 0.1~$\mu$L/min, and the needle was left in place for 5~minutes after injection to minimize backflow. Mice were allowed a minimum of two weeks for viral expression and surgical recovery.

\subsection{Contextual Fear Conditioning and Engram Labeling}

Following recovery, mice underwent contextual fear conditioning (CFC) to create the negative memory. The CFC protocol consisted of placement in Context A for 300 seconds, during which four foot shocks (1.5~mA, 2-second duration each) were delivered. Immediately following CFC, mice received an intraperitoneal injection of 4-hydroxytamoxifen (4-OHT) to permanently tag active neurons via TRAP2-mediated Cre-dependent recombination.

\subsection{Chronic Chemogenetic Activation}

Beginning two weeks after engram labeling, mice received chronic intraperitoneal injections of clozapine N-oxide (CNO) to reactivate the tagged negative-memory ensemble. Injections were administered over a three-month period. Two age cohorts were established: a young cohort (surgery at 3 months, testing at 6 months) and an old cohort (surgery at 11 months, testing at 14 months).

\subsection{Behavioral Testing}

After the three-month stimulation period, all mice underwent a comprehensive behavioral battery in the following order:

\textbf{Open field test (OFT):} Mice were placed in a 40~$\times$~40~cm arena for 10~minutes. Total distance traveled, time in center zone, center entries, and average speed were quantified using automated tracking software.

\textbf{Zero maze:} Mice were placed on an elevated zero maze for 5~minutes. Time in open arms, open arm entries, total distance, and average speed were recorded.

\textbf{Y-maze:} Spontaneous alternation was assessed in a Y-shaped maze with three equal arms. Mice freely explored for 8~minutes, and the percentage of spontaneous alternations and total arm entries were calculated.

\textbf{Fear extinction:} Mice were re-exposed to Context A (the CFC context) across multiple extinction sessions without foot shocks. Freezing behavior was quantified as the percentage of time spent immobile.

\textbf{Fear generalization:} Mice were placed in a novel Context B that differed from Context A in visual, tactile, and olfactory features. Freezing in the novel context was measured to assess fear generalization.

\subsection{RNA-Sequencing and Gene Set Enrichment Analysis}

A separate cohort of C57BL/6J mice received bilateral vHPC injections of a 1:1 cocktail of AAV9-c-Fos-tTA and AAV9-TRE-eYFP. Mice were maintained on doxycycline (DOX) diet to suppress tTA-TRE binding. DOX was removed on day 10 post-surgery to open the tagging window. On day 11, mice underwent either CFC (negative experience) or male-to-female social interaction (positive experience). DOX was immediately reinstated after the behavioral experience. Mice were sacrificed 24 hours later, and eYFP$^+$ engram cells were isolated by fluorescence-activated cell sorting (FACS). RNA was extracted and sequenced, followed by gene set enrichment analysis (GSEA) using curated pro-inflammatory and anti-inflammatory gene sets.

\subsection{Immunohistochemistry and Morphological Analysis}

Following behavioral testing, mice were transcardially perfused with 4\% paraformaldehyde. Brains were sectioned at 40~$\mu$m and processed for immunohistochemistry using antibodies against Iba1 (microglia), GFAP (astrocytes), and NeuN (neurons). Microglial and astrocytic morphology was quantified using 3DMorph, a MATLAB-based automated analysis pipeline \citep{york2018microglia}. Parameters assessed included ramification index, cell territory volume, average centroid distance, number of endpoints, average branch length, number of branch points, and longest branch length. NeuN$^+$ cell counts were quantified in the subiculum, ventral dentate gyrus (vDG), and vCA1 to assess neuronal survival.

\subsection{Statistical Analysis}

All behavioral and immunohistochemical data were analyzed using two-way ANOVA with treatment (hM3Dq versus mCherry) and age (6-month versus 14-month) as factors, followed by Tukey's post-hoc multiple comparisons test. Repeated-measures two-way ANOVA with Tukey's post-hoc test was used for CFC acquisition and extinction data. Statistical significance was set at $p < 0.05$. All analyses were performed blinded to treatment condition.

\section{Results}

\subsection{Negative vHPC Engram Cells Express Pro-Inflammatory Gene Signatures}

To characterize the molecular identity of positive and negative vHPC engram cells, we performed RNA-sequencing on FACS-sorted eYFP$^+$ cells from mice that had undergone either contextual fear conditioning (negative) or social interaction (positive). GSEA revealed that negative engram cells were significantly enriched for pro-inflammatory gene expression compared to untagged cells, while positive engram cells showed enrichment for anti-inflammatory gene sets (Table~\ref{tab:gsea}).

\begin{table}[ht]
\centering
\caption{Gene Set Enrichment Analysis of vHPC Engram Cells}
\label{tab:gsea}
\begin{tabular}{lcc}
\toprule
\textbf{Gene Set} & \textbf{Negative Engram} & \textbf{Positive Engram} \\
\midrule
Anti-inflammatory genes & Lower vs. untagged & Enriched vs. untagged \\
Pro-inflammatory genes & Enriched vs. untagged & Lower vs. untagged \\
\bottomrule
\end{tabular}
\end{table}

Among the most prominently upregulated genes in negative engram cells were \textit{Spp1} (secreted phosphoprotein 1/osteopontin), an extracellular matrix protein known to mediate inflammatory responses and upregulated after chronic social stress; \textit{Ttr} (transthyretin), a serum transporter increased after acute and chronic stress; and \textit{C1qb} (complement C1q B chain), involved in complement-mediated inflammatory processes and oxidative stress (Table~\ref{tab:deg}).

\begin{table}[ht]
\centering
\caption{Notable Differentially Expressed Genes in Negative Engram Cells}
\label{tab:deg}
\begin{tabular}{llp{7cm}}
\toprule
\textbf{Gene} & \textbf{Full Name} & \textbf{Function} \\
\midrule
\textit{Spp1} & Secreted phosphoprotein 1 & Extracellular matrix protein; mediates inflammatory responses; upregulated after chronic social stress \\
\textit{Ttr} & Transthyretin & Serum transporter for thyroid hormones; increased after acute/chronic stress \\
\textit{C1qb} & Complement C1q B chain & Complement-mediated inflammatory processes and oxidative stress \\
\bottomrule
\end{tabular}
\end{table}

\subsection{Chronic Negative-Engram Activation Increases Anxiety-Like Behavior}

Following three months of chronic CNO-mediated reactivation of the negative engram, both young (6-month) and aged (14-month) hM3Dq mice exhibited significantly increased anxiety-like behavior compared to mCherry controls. In the open field test, hM3Dq mice in both age groups spent significantly less time in the center zone and made fewer center entries, while total distance traveled and average speed were unaffected (Table~\ref{tab:anxiety}). This pattern confirms heightened anxiety without locomotor confounds.

In the zero maze, hM3Dq mice similarly spent less time in the open arms and made fewer open arm entries, without changes in total distance or average speed. Two-way ANOVA revealed significant main effects of treatment for center time ($p < 0.05$) and open arm time ($p < 0.05$), with Tukey's post-hoc comparisons confirming significant differences between hM3Dq and mCherry groups at both ages.

\begin{table}[ht]
\centering
\caption{Anxiety-Related Behavioral Measures After Chronic Negative-Engram Activation}
\label{tab:anxiety}
\begin{tabular}{lcccc}
\toprule
\textbf{Measure} & \textbf{6mo mCherry} & \textbf{6mo hM3Dq} & \textbf{14mo mCherry} & \textbf{14mo hM3Dq} \\
\midrule
OFT distance & Baseline & No change & Baseline & No change \\
OFT center time & Baseline & Decreased$^{*}$ & Baseline & Decreased$^{*}$ \\
OFT center entries & Baseline & Decreased$^{*}$ & Baseline & Decreased$^{*}$ \\
OFT average speed & Baseline & No change & Baseline & No change \\
ZM distance & Baseline & No change & Baseline & No change \\
ZM open arm time & Baseline & Decreased$^{*}$ & Baseline & Decreased$^{*}$ \\
ZM open arm entries & Baseline & Decreased$^{*}$ & Baseline & Decreased$^{*}$ \\
ZM average speed & Baseline & No change & Baseline & No change \\
\bottomrule
\multicolumn{5}{l}{\small $^{*}p < 0.05$, Tukey's post-hoc following two-way ANOVA.}
\end{tabular}
\end{table}

\subsection{Age-Dependent Spatial Working Memory Impairment}

Assessment of spatial working memory in the Y-maze revealed an age-dependent vulnerability to chronic negative-engram activation. While 6-month-old hM3Dq mice showed spontaneous alternation rates comparable to mCherry controls (approximately 50\%), 14-month-old hM3Dq mice exhibited significantly reduced alternation, indicating impaired spatial working memory. Total arm entries did not differ between groups at either age, confirming that the deficit reflects a cognitive impairment rather than a locomotor difference. Two-way ANOVA revealed a significant treatment~$\times$~age interaction ($p < 0.05$), with Tukey's post-hoc tests confirming the impairment was specific to the 14-month hM3Dq group.

\subsection{Impaired Fear Extinction and Heightened Fear Generalization}

All groups showed comparable freezing during the initial CFC session, confirming equivalent fear acquisition prior to the chronic stimulation period. Analysis of fear extinction revealed an age-specific pattern: 6-month-old hM3Dq mice showed significantly impaired fear extinction, maintaining elevated freezing levels across extinction sessions compared to age-matched mCherry controls (repeated-measures two-way ANOVA with Tukey's post-hoc, $p < 0.05$). In contrast, 14-month-old hM3Dq mice did not show a statistically significant extinction impairment.

Fear generalization was assessed by measuring freezing in a novel context never paired with shock. Both young and old hM3Dq mice showed significantly heightened freezing in the novel context compared to their respective mCherry controls (two-way ANOVA with Tukey's post-hoc, $p < 0.05$), indicating fear generalization across both age groups. A summary of all behavioral effects organized by age and behavioral domain is presented in Table~\ref{tab:summary}.

\begin{table}[ht]
\centering
\caption{Summary of Behavioral Effects by Age and Treatment Group}
\label{tab:summary}
\begin{tabular}{lcc}
\toprule
\textbf{Behavioral Domain} & \textbf{6mo hM3Dq vs. mCherry} & \textbf{14mo hM3Dq vs. mCherry} \\
\midrule
Anxiety (OFT center time) & Increased anxiety & Increased anxiety \\
Anxiety (ZM open arms) & Increased anxiety & Increased anxiety \\
Working memory (Y-maze) & No impairment & Impaired \\
Fear extinction & Impaired & Not significantly impaired \\
Fear generalization & Heightened & Heightened \\
Locomotor activity & No change & No change \\
\bottomrule
\end{tabular}
\end{table}

\subsection{Microglial Alterations Following Chronic Negative-Engram Activation}

Immunohistochemical analysis of Iba1$^+$ microglia revealed significant changes following chronic negative-engram activation. The number of Iba1$^+$ cells was altered in hippocampal subregions of hM3Dq mice compared to mCherry controls (two-way ANOVA with Tukey's post-hoc, $p < 0.05$). Detailed morphological analysis using 3DMorph revealed that chronic stimulation significantly altered multiple microglial morphological parameters, including ramification index, cell territory volume, average centroid distance, number of endpoints, average branch length, number of branch points, and longest branch length (all $p < 0.05$, two-way ANOVA with Tukey's post-hoc). These changes are consistent with a shift toward a more activated, reactive microglial phenotype typically observed under chronic stress conditions.

\subsection{Astrocytic Alterations Parallel Microglial Changes}

GFAP$^+$ astrocyte analysis revealed parallel alterations. The number of GFAP$^+$ cells was significantly changed in hM3Dq mice, and morphological parameters---including ramification index, cell territory volume, average centroid distance, number of endpoints, average branch length, number of branch points, and longest branch length---were all significantly altered (two-way ANOVA with Tukey's post-hoc, $p < 0.05$). These astrocytic changes mirrored the pattern observed in microglia, suggesting a coordinated neuroinflammatory response to chronic negative-engram reactivation.

\subsection{No Neuronal Loss Following Chronic Stimulation}

To determine whether the observed behavioral and glial changes were accompanied by neuronal death, we quantified NeuN$^+$ cells in the subiculum, ventral dentate gyrus, and vCA1. Two-way ANOVA revealed no significant differences between hM3Dq and mCherry groups at either age in any hippocampal subregion (all $p > 0.05$; $n = 3$ mice per group, 18 single-tile ROIs per brain region). This finding indicates that chronic chemogenetic reactivation of the negative engram did not induce detectable neuronal death, and that the observed behavioral and glial abnormalities reflect functional rather than degenerative changes.

\section{Discussion}

The present study demonstrates that chronic reactivation of a negative memory engram in the ventral hippocampus is sufficient to produce lasting behavioral and cellular abnormalities that closely parallel the effects of chronic external stress. This finding has important implications for understanding how rumination---the repetitive internal reactivation of negative memories---may contribute to the development and maintenance of neuropsychiatric disorders.

Our RNA-sequencing data reveal a fundamental molecular distinction between positive and negative vHPC engram cells: negative engrams are enriched for pro-inflammatory gene signatures, including \textit{Spp1}, \textit{Ttr}, and \textit{C1qb}, while positive engrams show anti-inflammatory enrichment. This molecular asymmetry suggests that the emotional valence of a memory is encoded not only in the pattern of neuronal activation but also in the inflammatory signaling potential of the engram cells themselves. The upregulation of complement component \textit{C1qb} in negative engram cells is particularly noteworthy, as complement signaling has been implicated in synaptic pruning and neuroinflammation.

The behavioral consequences of chronic negative-engram activation were both robust and multifaceted. Anxiety-like behavior was consistently elevated across both ages and two independent tests (open field and zero maze), while locomotor activity was preserved, ruling out motor confounds. Fear generalization was also observed in both age groups, consistent with the clinical observation that chronic stress promotes overgeneralization of threat-related memories. Importantly, however, we found striking age-dependent dissociations: spatial working memory impairment was observed only in 14-month-old mice, while fear extinction deficits were specific to 6-month-old mice. These dissociations suggest that the aging brain may be particularly vulnerable to the cognitive consequences of chronic negative-memory reactivation, while the younger brain may be more susceptible to disruptions in fear extinction circuits.

The neuroinflammatory changes observed in both microglia and astrocytes provide a cellular basis for the behavioral abnormalities. The morphological alterations in Iba1$^+$ microglia---including changes in ramification index, territory volume, and branch complexity---are consistent with a shift from a homeostatic surveillance state to a reactive, inflammatory phenotype \citep{nimmerjahn2005resting, york2021microglial}. Parallel changes in GFAP$^+$ astrocytes suggest a coordinated glial response. Crucially, NeuN$^+$ cell counts confirmed that these changes occurred without neuronal loss, indicating that the pathology is functional rather than neurodegenerative in nature. This pattern mirrors findings in chronic stress models, where behavioral impairments often precede and occur independently of neuronal death.

Our findings extend the engram literature in several important ways. While previous studies demonstrated that acute engram reactivation can elicit specific behavioral responses \citep{ramirez2013creating, redondo2014bidirectional}, our work shows that chronic reactivation produces a fundamentally different outcome: a broad, stress-like syndrome affecting multiple behavioral domains and accompanied by lasting glial changes. This suggests that the temporal pattern of engram reactivation---not just the content of the memory---is a critical determinant of its pathological consequences.

Several limitations should be acknowledged. First, while our study demonstrates that chronic negative-engram activation is sufficient to produce stress-like pathology, we did not directly compare this to chronic positive-engram activation, which would provide an important control for the effects of nonspecific chronic neuronal activation. Second, the use of CNO as the DREADD activator raises the possibility of off-target effects, although the mCherry control group addresses this concern. Third, our study was limited to male mice, and future work should examine sex-dependent effects.

In conclusion, our study provides the first evidence that chronic reactivation of a single negative memory engram can recapitulate the multi-domain behavioral and cellular pathology typically produced by chronic external stress. These findings suggest a neurobiological mechanism by which rumination and persistent negative memory retrieval could contribute to the development of anxiety, cognitive decline, and neuroinflammation, with age-dependent vulnerability profiles that may inform understanding of late-life neuropsychiatric risk.

\section{Conclusion}

This study establishes that the chronic reactivation of a negative memory engram in the ventral hippocampus, achieved through three months of chemogenetic stimulation in TRAP2 mice, produces lasting anxiety-like behavior, impaired fear extinction, fear generalization, age-dependent working memory deficits, and neuroinflammatory glial changes---all without neuronal loss. RNA-sequencing revealed that negative engram cells are intrinsically enriched for pro-inflammatory gene expression, providing a molecular basis for the inflammatory consequences of chronic reactivation. These results demonstrate that the persistent internal reactivation of a negative memory is sufficient to produce a chronic-stress-like phenotype, offering new insight into the neurobiological mechanisms underlying rumination-associated psychiatric disorders and highlighting the ventral hippocampal negative engram as a potential therapeutic target.

\bibliographystyle{plainnat}
\bibliography{references}

\end{document}
